\section*{Tóm tắt}
\thispagestyle{empty}
Sự phát triển vượt bậc của trí tuệ nhân tạo và xử lý ngôn ngữ tự nhiên đã mở đường cho sự ra đời giải pháp chatbot thông minh, hỗ trợ doanh nghiệp tự động hóa dịch vụ chăm sóc khách hàng. Đa số doanh nghiệp, từ trước đến nay, xây dựng chatbot dựa trên các mô hình học sâu được huấn luyện và triển khai cục bộ. Điều này dẫn đến khó khăn trong quá trình tái huấn luyện khi tri thức thay đổi. 

Chính vì vậy, đồ án chuyên ngành này tập trung vào việc thiết kế một nền tảng web để cung cấp dịch vụ chatbot dựa trên đám mây. Nền tảng "Chatbot as a service" này cho phép doanh nghiệp dễ dàng tích hợp chatbot tùy chỉnh vào website mà không cần chuyên môn kỹ thuật sâu rộng. Hệ thống đề xuất sẽ tận dụng sức mạnh của các mô hình ngôn ngữ lớn, cơ sở dữ liệu vector đám mây, và khả năng truy xuất tri thức hiệu quả từ kỹ thuật tăng cường truy xuất nâng cao (Advanced RAG) để cung cấp phản hồi chất lượng cao.

Kết quả đồ án cho thấy, nền tảng được thiết kế thể hiện được các tính năng nổi trội bao gồm giao diện thân thiện để cấu hình chatbot, chỉnh sửa tri thức, đào tạo chatbot và tích hợp qua API. Ngoài ra, hệ thống còn hỗ trợ cả trò chuyện tự động và nhân viên tư vấn trực tiếp, cùng với quản lý lịch hẹn (cho nhân viên) và đặt lịch hẹn (cho khách hàng) trên một giao diện thống nhất. 

Thông qua quá trình phân tích đánh giá về mặt lý thuyết, các công nghệ và mô hình sau đã được lựa chọn để thực hiện dự án: Frontend (ReactJS, Vite, NextUI, TailwindCSS), Backend (Django, Postgresql), AI (Gemini-1.5-flash model, LangChain, Pinecone). Hệ thống được đánh giá cao về tính linh hoạt, khả năng mở rộng dữ liệu và là cơ sở cho các cải tiến trong tương lai, hướng tới một giải pháp chatbot thương mại hóa. \\

\textbf{Từ khóa:} Artificial Intelligent, Natural Language Processing, Deep Neural Network, On-premise chatbot, Web-based platform, Cloud-based chatbot, Chatbot as a service, Large Language Models, Cloud-based Vector Database, Advanced RAG.

\clearpage
\pagenumbering{arabic}
